\section{Objeto de estudio}
\label{sec:objeto-estudio}

Esta sección corresponde a la \textbf{situación subyacente} de la investigación: abstracción y
caracterización del objeto de estudio antes del problema tecnológico (§1.3).

El \textbf{objeto de estudio} son los \textbf{contratos OpenAPI} (OpenAPI~3.x) con los que el
\textbf{sector bancario} describe sus APIs: artefactos documentales observables y medibles,
contrastables con modelos de negocio de referencia.

Quedan fuera del objeto la organización BIAN en su totalidad, el código de \textit{backends} y la
infraestructura de despliegue. BIAN entra como \textbf{modelo de referencia} (\ServiceDomain{}:
\textit{behaviors}, \textit{business objects}, relaciones de control) contra el cual se leen los
contratos.

BIAN también publica contratos OpenAPI de referencia, pero el estudio no reduce el estándar a su
capa REST: el \ServiceDomain{} expresa el negocio y el \textbf{contrato OpenAPI bancario}
---del banco, de un integrador o de un catálogo--- es la \textbf{capa técnica evaluada}. Esa
distinción orienta la lectura del contrato frente al \ServiceDomain{} y se ilustra en Anexo~F.

Para evitar ambigüedades, se distinguen tres niveles. La \Cref{fig:ontologia} resume la relación
entre objeto, referencia BIAN, unidad de análisis y producto del modelo; el detalle tabular sigue
en la \Cref{tab:objeto-unidad-producto}.

% Figura ontología conceptual — elaboración propia Grupo 9.2 (TikZ)
\begin{figure}[H]
  \centering
  \begin{tikzpicture}[
    font=\small\sffamily,
    box/.style={draw, rounded corners=2pt, align=center, minimum width=2.6cm,
                minimum height=0.85cm, fill=gray!8},
    ref/.style={box, fill=blue!8},
    prod/.style={box, fill=green!10},
    arr/.style={-{Stealth[length=2.2mm]}, thick, gray!70}
  ]
    \node[box] (sector) {Sector bancario\\APIs documentadas};
    \node[box, below=0.9cm of sector] (obj) {\textbf{Objeto:}\\contrato OpenAPI bancario};
    \node[ref, right=1.1cm of obj] (bian) {\textbf{Referencia:}\\\ServiceDomain{} BIAN};
    \node[box, below=0.9cm of obj] (ua) {\textbf{UA:}\\endpoint/schema $\leftrightarrow$\\ \Behavior{} / \BusinessObject{}};
    \node[prod, below=0.9cm of ua] (out) {\textbf{Producto:}\\S, E, C, \AlignmentScore{}};

    \draw[arr] (sector) -- (obj);
    \draw[arr] (obj) -- (ua);
    \draw[arr] (bian) -- (ua);
    \draw[arr] (ua) -- (out);
  \end{tikzpicture}
  \caption{Mapa conceptual del estudio: objeto, referencia, unidad de análisis y producto del
           modelo (elaboración propia).}
  \label{fig:ontologia}
\end{figure}


\begin{table}[H]
  \centering
  \caption{Objeto, unidad de análisis y producto del modelo}
  \label{tab:objeto-unidad-producto}
  \begin{tabular}{@{}p{0.28\textwidth}p{0.62\textwidth}@{}}
    \toprule
    \textbf{Concepto} & \textbf{En este estudio} \\
    \midrule
    \textbf{Objeto de estudio} & El contrato OpenAPI bancario completo (artefacto documental) \\
    \textbf{Unidad de análisis} &
    Elementos comparables dentro del \textit{pipeline}: endpoint/operation/schema $\leftrightarrow$
    \Behavior{} / \BusinessObject{} BIAN \\
    \textbf{Producto del modelo} & Scores S, E, C y \AlignmentScore{} por contrato evaluado \\
    \bottomrule
  \end{tabular}
\end{table}

En el procedimiento propuesto, el \textbf{objeto} es el archivo de entrada (el contrato); la
\textbf{unidad de análisis} es cada par procesado dentro de él (endpoint--
\Behavior{}, esquema--\BusinessObject{}, etc.); el \textbf{producto} agrega esos resultados en
\AlignmentScore{} y clasificación por instancia \texttt{contrato} $\times$ \ServiceDomain{}.

Cada instancia se identifica de forma unívoca, p. ej.
\texttt{\{fuente\}\_\{service-domain\}\_\{version\}}, coherente con la adquisición de datos del
Capítulo~\ref{cap:metodologia} (§4.3).

\begin{table}[H]
  \centering
  \caption{Ejemplo genérico ilustrativo (sin caso piloto fijo)}
  \label{tab:ejemplo-generico}
  \begin{tabular}{@{}p{0.44\textwidth}p{0.44\textwidth}@{}}
    \toprule
    \textbf{Elemento OpenAPI (genérico)} & \textbf{Elemento BIAN de referencia (ilustrativo)} \\
    \midrule
    \texttt{POST /payment-orders} & \Behavior{} \texttt{Initiate} del dominio de pagos \\
    \texttt{GET /payment-orders/\{id\}} & \Behavior{} \texttt{Retrieve} \\
    Esquema \texttt{PaymentOrder} (propiedades \texttt{amount}, \texttt{currency}) &
    \BusinessObject{} \texttt{Payment} / \texttt{PaymentAmount} \\
    \bottomrule
  \end{tabular}
\end{table}


La tabla es \textbf{explicativa}: nombres ficticios para mostrar la unidad de análisis. El caso
ampliado de \textit{Payment Initiation} está en Anexo~F; el procedimiento aplica a cualquier
contrato OpenAPI bancario emparejado con un \ServiceDomain{} BIAN.
La \Cref{fig:objeto-unidad} resume objeto, unidad de análisis, producto del modelo y el ejemplo
genérico.

\begin{figure}[H]
  \centering
  \includegraphics[width=0.92\textwidth]{Figura_02_Objeto_Unidad_Analisis}
  \caption{Objeto de estudio, unidad de análisis y producto (ejemplo genérico)}
  \label{fig:objeto-unidad}
\end{figure}

\subsection{Producto central: modelo de alineación}

El \textbf{resultado esperado} ---lo que la matriz y el Capítulo~\ref{cap:objetivos} declaran como
producto--- es un \textbf{modelo de alineación}. Dado un contrato OpenAPI bancario y un
\ServiceDomain{} BIAN emparejado, debe calcular:

\begin{table}[H]
  \centering
  \caption{Componentes del modelo de alineación}
  \label{tab:componentes-score}
  \begin{tabular}{@{}p{0.28\textwidth}p{0.62\textwidth}@{}}
    \toprule
    \textbf{Componente} & \textbf{Significado} \\
    \midrule
    \SemanticScore{} & Similitud de significado entre conceptos OpenAPI y BIAN \\
    \StructuralScore{} & Correspondencia de jerarquía y organización (REST $\leftrightarrow$ SD) \\
    \CoverageScore{} & Proporción de conceptos BIAN cubiertos por el contrato \\
    \bottomrule
  \end{tabular}
\end{table}

A partir de estos componentes se obtiene el \AlignmentScore{}:

\begin{quote}
  \AlignmentScore{} $= \alpha \cdot S + \beta \cdot E + \gamma \cdot C$, con $\alpha + \beta +
  \gamma = 1$
\end{quote}

La clasificación interpretativa del \textit{score} es: \textbf{Alta} ($\geq 0{,}80$),
\textbf{Media} (0,60--0,79), \textbf{Baja} (0,40--0,59) y \textbf{Nula} ($< 0{,}40$). Los umbrales
son \textbf{convención provisional del diseño} para este plan; su calibración empírica
permanece fuera del alcance (§1.6). Los scores S, E y C se calculan por unidad de análisis (par UA); el
\AlignmentScore{} y la clasificación se reportan por instancia contrato $\times$ \ServiceDomain{}.
La \Cref{fig:alignment-score} ilustra el cálculo del score de alineación.

\begin{figure}[H]
  \centering
  \figCapThree
  \caption{Cálculo del score de alineación: S, E, C $\rightarrow$ \AlignmentScore{}}
  \label{fig:alignment-score}
\end{figure}

Esas cinco fases de transformación ---diseño, normalización, cálculo de E, S y C, integración en
\AlignmentScore{}--- coinciden con la matriz de consistencia (Anexo~C) y con la cadena OE1--OE4 del
Capítulo~\ref{cap:objetivos}.
